\documentclass[../main]{subfiles}
\begin{document}
\chapter{Motor de Carnot}
\img{images/07/carnot.jpg}{Motor ideal de Carnot. }
\info{
Contenido}{
  Ciclo Carnot, rendimiento, diagramas, ciclos reales.
}
\info{Ciclo de Carnot}{
  El ciclo de Carnot es un ciclo termodinámico que se produce en un
  equipo o máquina cuando trabaja absorbiendo una cantidad de calor
  $Q_1$ de una fuente de mayor temperatura y cediendo un calor $Q_2$
  a la de menor temperatura produciendo un trabajo sobre el exterior.
}
\section{Rendimiento}
El rendimiento $\eta$ de este ciclo viene definido por
\img{images/07/carnotciclo.png}{Esquema de una máquina de Carnot.}

\begin{equation}
  \eta = \frac{W}{Q_1}
\end{equation}

Como

$Q_1=Q_2+W \Rightarrow W=Q_1-Q_2$

Entonces

\begin{equation}
  \eta = \frac{Q1-Q_2}{Q_1}=1-\frac{Q_2}{Q_1}
\end{equation}

\begin{equation}
  \eta =1-{\frac {T_{2}}{T_{1}}}
\end{equation}

Siendo:

$W$: Trabajo $[kcal,Joule,kgm]$

$\eta$: Rendimiento valor entre 1 y 0.

$T_2,T_1$: Temperaturas de foco frío y caliente en grados Kelvin.

$Q_1,Q_2$: Calor de entrada y salida.$[kcal,Kg m,Joule]$

\begin{equation}
  \eta =1-{\frac {T_{2}}{T_{1}}}=1-{\frac {Q_{2}}{Q_{1}}}
\end{equation}
\begin{equation}
  \eta_{\%} =\left( 1-\frac {T_{2}}{T_{1}}\right)\times100
\end{equation}
Una máquina térmica que realiza este ciclo se denomina máquina de Carnot.
Como todos los procesos que tienen lugar en el ciclo ideal son
reversibles, el ciclo puede invertirse y la máquina absorbería calor
de la fuente fría y cedería calor a la fuente caliente, teniendo que
suministrar trabajo a la máquina. Si el objetivo de esta máquina es
extraer calor de la fuente fría (para mantenerla fría) se denomina
máquina frigorífica, y si es ceder calor a la fuente caliente, bomba de calor.
\actividad{Contestar el cuestionario}
\begin{enumerate}
  \item Explicar el ciclo de Carnot: Etapas, y realizar diagrama P-V
    \emph{(presión-volumen)} y T-S \emph{(temperatura-entropía)}.
  \item Explique el primer y segundo teorema de Carnot.
  \item ¿Cómo se calcula el rendimiento de una máquina de Carnot?
  \item ¿Qué diferencia hay entre los ciclos reales y el ciclo de Carnot?
\end{enumerate}
\actividad{Problemas}
\begin{enumerate}
  \item Supóngase que se tiene una máquina térmica capaz de funcionar
    de los siguientes modos:
    \begin{enumerate}
      \item Entre \qty{200}{\degree\kelvin}  y \SI{400}{\degree\kelvin},
      \item Entre \SI{600}{\degree\kelvin} y \SI{400}{\degree\kelvin}.
      \item Entre \SI{100}{\celsius} y \SI{20}{\celsius}
      \item Entre \SI{473}{\celsius} y \SI{683}{\celsius}
    \end{enumerate}
    ¿Cuál es rendimiento del ciclo del ciclo ideal de Carnot que
    podría tener en cada caso?
  \item Una máquina térmica con eficiencia de $\eta_{\%}=22\%$
    produce $W=1530J$ de trabajo. Encontrar:
    \begin{enumerate}
      \item La cantidad de calor absorbida del depósito térmico.
      \item La cantidad de calor desechado al depósito térmico.
    \end{enumerate}
  \item La máquina térmica con eficiencia de $\eta_{\%}=87\%$ produce
    $W=4186J$ de trabajo. Encontrar:
    \begin{enumerate}
      \item La cantidad de calor absorbida del depósito térmico.
      \item La cantidad de calor desechado al depósito térmico.
    \end{enumerate}
  \item Una máquina según un ciclo de Carnot reversible absorbiendo
    calor de un tanque de agua a $\SI{100}{\celsius}$ y cediéndolo al
    aire en el interior de un local que se mantiene a
    $\SI{26}{\celsius}$ La máquina produce un trabajo de $41860J$. Calcule:
    \begin{enumerate}
      \item El rendimiento $\eta$ de la máquina de Carnot.
      \item La cantidad de calor cedido por la fuente caliente $Q_1$.
      \item la cantidad de calor entregado en la fuente fría $Q_2$.
    \end{enumerate}
  \item Una máquina de ciclo ideal tiene un rendimiento pobre del
    $\eta_{\%}=10\%$ Si produce un trabajo de W=1000J, estando la
    fuente fría a $T_2=\SI{10}{\celsius}$. Calcular los calores
    entregados y recibidos $Q_1,Q_2$, la temperatura de la fuente
    caliente $T_1$.
  \item Para el funcionamiento de una máquina, se pierde el 30\% del
    calor entregado, si la fuente caliente está a
    $T_1=\SI{400}{\celsius}$ y se realiza un trabajo de 200J Calcular:
    \begin{enumerate}
      \item El calor perdido $Q_2$.
      \item El rendimiento $\eta$.
      \item El calor entregad $Q_1$.
      \item calcular $T_1$ y $T_2$
      \item Si se quiere duplicar el trabajo realizado ¿Cuál debería
        ser la temperatura de la fuente caliente $T_1$?
    \end{enumerate}
  \item Un vehículo de $P=2300kg$ que funciona mediante el ciclo de
    Carnot con un rendimiento de $\eta=95\%$, necesita subir una
    cuesta de $h=10m$ de altura, la fuente fría es la temperatura del
    ambiente de $T_2=\SI{30}{\celsius}$ ¿Qué temperatura debe tener
    la fuente caliente para que pueda subir a esa altura?
  \item Dos ciclos de Carnot funciona en serie, end dos etapas, si el
    primer ciclo si las temperaturas correspondientes son
    $\SI{10}{\celsius}$ $\SI{150}{\celsius}$ y $\SI{400}{\celsius}$.
    Calcular el rendimiento en cada una de las etapas, y el trabajo
    en cada una si el calor entregado inicialmente es de $100J$.
  \item Una máquina de Carnot trabaja entre $\SI{100}{\celsius}$ y
    $\SI{1000}{\celsius}$. Calcular el  cantidad de calor necesaria
    para producir el trabajo de 1J.
  \item Dos máquinas de Carnot se conectan en paralelo, entre
    $\SI{100}{\celsius}$ y $\SI{500}{\celsius}$, la primer máquina
    produce un trabajo de 100J y la segunda 200J. Calcular el
    rendimiento de cada una y del conjunto.
  \item Siendo la temperatura de la fuente fría $\SI{100}{\celsius}$
    y el rendimiento $\eta=75\%$, determinar la temperatura de la
    fuente caliente.
  \item Siendo la temperatura de la fuente fría $\SI{200}{\celsius}$
    y el rendimiento $\eta=55\%$, determinar la temperatura de la
    fuente caliente.
  \item Al fluido que evoluciona dentro de una máquina térmica ideal
    de CARNOT se le suministra una cantidad de calor igual a
    $100000kcal$. Suponiendo que el foco caliente se encuentra a
    $\SI{127}{\celsius}$ y el foco frío a $\SI{7}{\celsius}$, calcular:
    \begin{enumerate}
      \item el rendimiento térmico de la máquina.
      \item el trabajo desarrollado; y
      \item la cantidad de calor que pasa al foco frío.
    \end{enumerate}
  \item Una máquina térmica que funciona con ciclo ideal de CARNOT y
    cuyo rendimiento térmico es de $\eta\%=30\%$, tiene el foco frío
    a$\SI{7}{\celsius}$. Calcular la temperatura del foco caliente.
  \item Calcular la eficiencia de una máquina frigorífica que
    funciona con el ciclo ideal de CARNOT, si el foco caliente se
    encuentra a temperatura de $\SI{127}{\celsius}$ y el foco frío a
    $\SI{7}{\celsius}$
  \item Un fluido que evoluciona dentro de una máquina frigorífica
    con ciclo ideal de CARNOT entrega en el foco caliente $100000
    kcal$. Suponiendo que el foco frío se encuentra a
    $\SI{7}{\celsius}$ y el foco caliente en $\SI{127}{\celsius}$. Calcular:
    \begin{enumerate}
      \item la cantidad de calor extraída del foco frío;
      \item el trabajo mecánico que es necesario suministrar; y
      \item la eficiencia de la máquina $\eta$.
    \end{enumerate}
\end{enumerate}
\end{document}
